% Generated draft: requires mathematical and source review.
\documentclass[11pt]{article}
\usepackage[margin=25mm]{geometry}
\usepackage{fontspec}
\setmainfont{Latin Modern Roman}
\usepackage{amsmath,amssymb,mathtools,unicode-math}
\setmathfont{Latin Modern Math}
\usepackage{enumitem,needspace,xurl,hyperref}
\setlength{\parindent}{0pt}
\setlength{\parskip}{5pt}
\setlength{\emergencystretch}{3em}
\newcommand{\N}{\mathbb{N}}
\newcommand{\Z}{\mathbb{Z}}
\newcommand{\Q}{\mathbb{Q}}
\newcommand{\R}{\mathbb{R}}
\newcommand{\C}{\mathbb{C}}
\newcommand{\F}{\mathbb{F}}
\newcommand{\A}{\mathbb{A}}
\newcommand{\PP}{\mathbb{P}}
\newcommand{\E}{\mathbb{E}}
\newcommand{\eps}{\varepsilon}
\newcommand{\dd}{\,\mathrm{d}}
\begin{document}
\section*{UnsolvedMath ID: 8800042}
\textbf{Source problem code:} catalogue entry
\subsection{Definitions and mathematical statement}
For a prime power \(q\) and an elliptic curve \(E/\mathbb F_q\), put \(N_q(E)=\#E(\mathbb F_q)\). For an integer \(N>1\), let \(P^+(N)\) be its largest prime divisor, with \(P^+(1)=1\). The source asks to quantify curves for which \(P^+(N_q(E))\) is large.

For example, if \(\mathcal E_q\) is the finite set of \(\mathbb F_q\)-isomorphism classes, one possible counting function is
\[
 A(q,Y)=\#\{[E]\in\mathcal E_q:P^+(N_q(E))\ge Y\},\qquad Y>1.
\]
Determine useful sharp bounds or asymptotics, or the corresponding proportion, for the intended curve-counting model and threshold \(Y=Y(q)\). This is an explicit possible formalization: the source does not choose isomorphism-class counting rather than Weierstrass-model or automorphism-weighted counting, and does not define ``large''. Thresholds such as \(q^\theta\) or proximity to \(q\) give different questions. Nor is the field-asymptotic regime specified.

The Hasse bound places \(N_q(E)\) in \([q+1-2\sqrt q,q+1+2\sqrt q]\), but the distribution of curve orders within this interval is part of the counting problem. The supplied summary mentions Lenstra-type distribution estimates and Ford--Shparlinski-type work on large prime factors, reporting partial quantitative results and some near-\(q\) prime-divisor claims. It provides no exact thresholds, proportions, hypotheses or references, so no universal near-\(q\) guarantee or optimal exponent is inferred. The 2006 framing and later partial-progress assessment remain source reports pending current primary-literature review.
\subsection{Short English statement}
Count elliptic curves whose finite-field group order has a large prime factor. A precise result must specify the curve-sampling measure, prime-factor threshold and field-asymptotic regime; the source gives only a broad quantitative research target.
\subsection{Sources}
\begin{itemize}
\item Local registry: \url{data/unsolvedmath/problems.json}, permanent UnsolvedMath ID \(8800042\), source code catalogue entry; supplied statement, background, and research summary, updated 2026-08-06. Attributed to Igor Shparlinski.
\item IPAM, \emph{Workshop on Number Theory and Cryptography---Open Problems} (2006), attachment 6224, PDF page 21: \url{https://www.ipam.ucla.edu/programs/workshops/workshop-i-number-theory-and-cryptography-open-problems/?tab=schedule}. Supplied PDF/legacy-PPT visual audit accessed 2026-07-29; it reports 19 inspected surviving attachments.
\item Lenstra and Ford--Shparlinski-type estimates are mentioned by the summary without full publication details or exact quantitative statements.
\end{itemize}
\end{document}
